\subsection{Examples of sheaves}

Let B be a topological space.

\subitemhead{Sheaves of mappings}

Let X be a set. For every open set U of B, denote by \(\mathcal{F}(\mathrm{U};\mathrm{X})\) the set of mappings from U into X (E, II, p.~31). For every pair (U, V) of open sets of B such that \(\mathrm{U}\subset \mathrm{V}\) , let \(r_{\mathrm{UV}}\colon \mathcal{F}(\mathrm{V};\mathrm{X})\to\mathcal{F}(\mathrm{U};\mathrm{X})\) be the restriction map \(f\mapsto f|\mathrm{U}\) . It is clear that the pair \((\mathcal{F}(\mathrm{U};\mathrm{X}),r_{\mathrm{UV}})\) is a sheaf on B. It is called the sheaf on B of mappings with values in X, and is denoted by \(\underline{\mathcal{F}} (\mathrm{B};\mathrm{X})\) .

\subitemhead{Sheaves of continuous mappings}

Let X be a topological space. For every open set U of B, let \(\mathcal{C}(\mathrm{U};\mathrm{X})\) be the set of continuous mappings from U into X. Then, \((\mathcal{C}(\mathrm{U};\mathrm{X}))\) is a subsheaf of the sheaf \(\underline{\mathcal{F}} (\mathrm{B};\mathrm{X})\) by prop. 4 of TG, I, p.~19. The sheaf thus obtained is denoted \(\underline{\mathcal{C}} (\mathrm{B};\mathrm{X})\) and called the sheaf on B of continuous mappings with values in X. In the particular case where X is equipped with the discrete topology, the sheaf \(\underline{\mathcal{C}} (\mathrm{B};\mathrm{X})\) takes the name of the sheaf on B of locally constant functions with values in X.

\subitemhead{Sheaves of continuous sections}

Let E be a topological space and let \(p \colon \mathrm{E} \to \mathrm{B}\) be a continuous map. For every open set U in B, we denote by \(\mathcal{S}(\mathrm{U};p)\) (or \(\mathcal{S}(\mathrm{U};\mathrm{E})\) when there is no possible confusion) the set of continuous sections of \(p\) over U. The family \((\mathcal{S}(\mathrm{U};p))\) is a subsheaf of the sheaf \(\underline{\mathcal{C}}(\mathrm{B};\mathrm{E})\). The sheaf thus obtained is denoted \(\underline{\mathcal{L}}(\mathrm{B};\mathrm{E})\) or simply \(\underline{\mathcal{S}}(\mathrm{E})\) and called the sheaf on B of continuous sections of the B-space \((\mathrm{E},p)\). We shall see in no. 6 below that every sheaf on B is isomorphic to the sheaf on B of continuous sections of an étale B-space.

\subitemhead{Sheaves of B-morphisms}

Let \((\mathrm{E},p)\) and \((\mathrm{E}^{\prime},p^{\prime})\) be B-spaces. For every open set U in E, we denote by \(\mathcal{C}_{\mathrm{B}}(\mathrm{U};\mathrm{E}^{\prime})\) the set of B-morphisms from \((\mathrm{U},p|\mathrm{U})\) into \((\mathrm{E}^{\prime},p^{\prime})\) . The family \((\mathcal{C}_{\mathrm{B}}(\mathrm{U};\mathrm{E}^{\prime}))\) is a subsheaf of the sheaf \(\underline{\mathcal{C}}(\mathrm{E};\mathrm{E}^{\prime})\). The sheaf thus obtained is denoted \(\underline{\mathcal{C}}_{\mathrm{B}}(\mathrm{E};\mathrm{E}^{\prime})\) and called the sheaf on E of B-morphisms with values in \((\mathrm{E}^{\prime},p^{\prime})\) . When \((\mathrm{E},p)\) is equal to \((\mathrm{B},\mathrm{Id}_{\mathrm{B}})\), this sheaf is the sheaf \(\underline{\mathcal{L}}(\mathrm{B};\mathrm{E})\) of Example 3.

For every open set U of B, let \(\mathcal{M}(\mathrm{U})\) be the set of U-morphisms from \(p^{-1}(\mathrm{U})\) into \((p')^{-1}(\mathrm{U})\) . For every pair \((\mathrm{U},\mathrm{V})\) of open subsets of B such that \(\mathrm{U}\subset \mathrm{V}\) , let \(m_{\mathrm{UV}}\colon \mathcal{M}(\mathrm{V})\to \mathcal{M}(\mathrm{U})\) be the map which assigns to a V-morphism \(f\colon p^{-1}(\mathrm{V})\to (p')^{-1}(\mathrm{V})\) the U-morphism of \(p^{-1}(\mathrm{U})\) into \((p')^{-1}(\mathrm{U})\)

deduced from \(f\) by passing to subsets. Then \((\mathcal{M}(\mathrm{U}), m_{\mathrm{UV}})\) is a sheaf on B. We denote it by \(\underline{\mathcal{M}or}_{\mathrm{B}}(\mathrm{E}; \mathrm{E}')\) and it is called the sheaf on B of B-morphisms of \((\mathrm{E}, p)\) into \((\mathrm{E}', p')\) .

For every open set U of B, let \(\mathcal{I}s(\mathrm{U})\) be the subset of \(\mathcal{M}(\mathrm{U})\) consisting of the U-isomorphisms of \(p^{-1}(\mathrm{U})\) into \((p')^{-1}(\mathrm{U})\) . The family \((\mathcal{I}s(\mathrm{U}))\) is a subsheaf of the sheaf \(\underline{\mathcal{Mor}}_{\mathrm{B}}(\mathrm{E};\mathrm{E}^{\prime})\) of morphisms of \((\mathrm{E},p)\) into \((\mathrm{E}^{\prime},p^{\prime})\). The sheaf thus obtained is denoted \(\underline{\mathcal{Isom}}_{\mathrm{B}}(\mathrm{E};\mathrm{E}^{\prime})\) and called the sheaf on B of B-isomorphisms of \((\mathrm{E},p)\) into \((\mathrm{E}^{\prime},p^{\prime})\) .

\subitemhead{Sheaves of mappings of class \(\mathbf{C}^r\)}

Let X and Y be varieties of class \(C^{r}\) over a field K (the conventions on K and r being those of VAR, R). For every open subset U of X, let \(\mathcal{C}^{r}(\mathrm{U};\mathrm{Y})\) be the set of mappings of class \(C^{r}\) from U into Y. The family \((\mathcal{C}^{r}(\mathrm{U};\mathrm{Y}))\) is a subsheaf of the sheaf \(\underline{\mathcal{C}}(\mathrm{X};\mathrm{Y})\). The sheaf thus obtained is denoted \(\underline{\mathcal{C}}^{r}(\mathrm{X};\mathrm{Y})\) and called the sheaf on X of mappings of class \(C^{r}\) with values in Y (cf. VAR, R, 5.4.2).

\subitemhead{Sheaves of subspaces}

If U and V are open subsets of B such that \(U \subset V\) , denote by \(i_{\mathrm{UV}}: \mathfrak{P}(V) \to \mathfrak{P}(U)\) the map which assigns to each subset A of V the subset \(A \cap U\) . The pair \((\mathfrak{P}(U), i_{\mathrm{UV}})\) is a sheaf, called the sheaf of subspaces of B and denoted \(\underline{\mathfrak{P}}(B)\) . Indeed, if one denotes by X the set \(\{0;1\}\), the map which assigns to every subset A of U its characteristic function \(\varphi_{A}^{U}: U \to X\) is a bijection from \(\mathfrak{P}(U)\) onto \(\mathcal{F}(U;X)\) (E, III, p.~38); moreover, if U and V are open sets such that \(U \subset V\) , for every subset A of V, \(\varphi_{A \cap U}^{U}\) is the restriction to U of \(\varphi_{A}^{V}\) so that \(\underline{\mathfrak{P}}(B)\) is identified with the sheaf on B of maps with values in X.

Let, for every open set U of B, \(\mathcal{L}(\mathrm{U})\) be a subset of \(\mathfrak{P}(\mathrm{U})\) . For \((\mathcal{L}(\mathrm{U}))\) to be a subsheaf of \(\underline{\mathfrak{P}}(\mathrm{B})\) it is necessary and sufficient that the following condition be satisfied:

(F') Let \((\mathrm{U}_{i})_{i\in\mathrm{I}}\) be a family of open sets of B, U their union, and let A be a

subset of U; for A to belong to \(\mathcal{L}(\mathrm{U})\) , it is necessary and

sufficient that, for every i in I, \(A \cap \mathrm{U}_{i}\) belongs to \(\mathcal{L}(\mathrm{U}_{i})\).

For example, if \(\mathcal{L}(\mathrm{U})\) is the set of closed subsets of U, the condition \((\mathrm{F}')\) is satisfied.

\subitemhead{Products of sheaves}

Let \(\mathcal{B}\) be a base of the topology of B and let I be a set. For every \(i \in I\) , let \(\mathcal{F}_{i} = (\mathcal{F}_{i}(U), f_{i,\mathrm{UV}})\) be a presheaf on B relative to the base \(\mathcal{B}\). For every open set \(U \in \mathcal{B}\) , set \(\mathcal{F}(U) = \prod_{i \in I} \mathcal{F}_{i}(U)\) , and for every pair (U, V) of elements of \(\mathcal{B}\) such that U ⊂ V, denote by \(f_{\mathrm{UV}}\) the map ( \(f_{i,\mathrm{UV}}\) ) \(_{i\in\mathrm{I}}\) : \(\mathcal{F}(\mathrm{V}) \to \mathcal{F}(\mathrm{U})\) . Then ( \(\mathcal{F}(\mathrm{U}), f_{\mathrm{UV}} \) ) is a presheaf on B relative to \(\mathcal{B}\); it is called the product presheaf of the family ( \(\mathcal{F}_i\) ) and denoted \(\prod_{i\in\mathrm{I}}\mathcal{F}_i\) . It is a sheaf if, for every \(i \in \mathrm{I}\) , \(\mathcal{F}_i\) is a sheaf.

